% =====================================================================================
\section{Discussion}
\label{sec:discussion}

\subsection{Ordering more than sizing}
\label{sec:ordering}

% Source: Section 3 numbers; measurement_ceiling_summary.json (coverage 0.8695, 1,149 units, 2024).
The model is better at saying which of two hitters will do better against a hand than at saying by how many points. Part of this is a property of the data rather than the model. Empirical Bayes sits just as close to the error floor (Section~\ref{sec:ceiling}), so one season of outcomes cannot show much skill at sizing a split for any method. The model's intervals are also too narrow for sizing, covering the observed value less often than their nominal rate (Section~\ref{sec:test}). The same limit applies to single hitters, whose observed splits are rarely precise enough to say which projection was closer. Many lineup decisions depend on order rather than size. Choosing which of two bench bats starts against a left-handed pitcher is a ranking question, and a few points of error in both projections does not change the call if the two are in the right order.

The model was built to be accurate for every hitter, with particular attention to hitters who have little history against a hand, and its advantage over the baselines varies with exposure. The strata cuts are a convention, so the pattern is better read as a curve than as three groups. With little exposure, the model has few pitches against that hand to place a hitter, and 43\% of low-stratum units belong to hitters absent from training, who receive one of four generic projections. The model's point estimates there still rank hitters better than both baselines, but the stratum is too small to confirm the margin. With heavy exposure, several seasons of outcomes already measure a hitter well, which leaves little for any method to add. In between, a hitter has seen enough pitches from a hand for the embedding to describe how the hitter handles them, while the outcomes from those plate appearances are still too few to stabilize. Pitch-level outcomes accumulate several times faster than plate appearance results (about four pitches per plate appearance), so this middle range is where they add the most over methods built on outcomes alone.

\subsection{Case studies}
\label{sec:anchors}

% Source: docs/phase-v-spec.md V.8 (anchors and expectations written before results);
% results/case_studies_2025/case_studies_2025.{csv,json}.
Five well-known hitters serve as case studies. Four had an expected platoon direction fixed before any projection was examined, and Alec Bohm replaced a pre-registered right-handed hitter whose training data showed no platoon split. Joc Pederson and Kyle Schwarber are left-handed hitters with reputations as platoon bats and should project worse against left-handed pitchers. Bohm, a right-handed hitter, should project better against them. Mike Trout and Juan Soto play every day and should project close to flat. All five are high-exposure hitters, so they show how the model describes hitters it knows well. Table~\ref{tab:cases} gives each projected 2025 split beside the empirical Bayes projection and the observed split.

\begin{table}[!ht]
\centering\small
\begin{tabular}{llrrrr}
\toprule
Hitter & Bats & Model & EB & Observed (SE) & $z$ \\
\midrule
Joc Pederson & L & $-0.047$ & $-0.039$ & $+0.172$ (0.121) & 1.77 \\
Kyle Schwarber & L & $-0.042$ & $-0.028$ & $+0.025$ (0.039) & 1.52 \\
Juan Soto & L & $-0.041$ & $-0.044$ & $-0.074$ (0.041) & 0.78 \\
Mike Trout & R & $+0.009$ & $+0.007$ & $-0.062$ (0.056) & 1.25 \\
Alec Bohm & R & $+0.022$ & $+0.033$ & $+0.042$ (0.052) & 0.28 \\
\midrule
Aaron Judge & R & $+0.020$ & $+0.003$ & $+0.038$ (0.049) & 0.54 \\
José Altuve & R & $+0.021$ & $+0.026$ & $+0.013$ (0.051) & 0.22 \\
Marcell Ozuna & R & $+0.020$ & $+0.002$ & $-0.006$ (0.050) & 0.33 \\
\bottomrule
\end{tabular}
\caption{2025 platoon split (wOBA against left-handed pitchers minus wOBA against right-handed pitchers) for the five case-study hitters (top) and three well-known high-exposure hitters among the twenty on whom the model and empirical Bayes disagreed most in 2024 (bottom). SE is the standard error of the observed split, and $z$ is the distance of the observed split from the midpoint of the two projections in standard errors. Every $z$ is below 1.96, so noise alone could flip which projection is closer.}
\label{tab:cases}
\end{table}

The projections match four of the five expectations. Pederson and Schwarber project worse against left-handed pitchers, Bohm better, and Trout close to flat. Soto, expected to be flat, projects 0.041 worse against left-handed pitchers, and the observed 2025 split, $-0.074$, runs the same way. Pederson's head outputs differ by pitcher hand. Pederson swings less than an average hitter and does more damage on balls in play against both hands, but the swing shortfall is larger and the damage edge slightly smaller against left-handed pitchers. In 2025, empirical Bayes was closer for all five, by 0.002 to 0.014, each well inside the noise (Figure~\ref{fig:cases}). Pederson had 19 PA against left-handed pitchers in 2025 and an observed split of $+0.172$ with a standard error of 0.121, wide enough to hold either projection and the opposite sign.


% Source: results/v_chain_clean/disagreement_summary.json, disagreement_top.csv;
% case_studies_2025.csv (disagreement rows, stratum by 2025 prior PA).
Hitters on whom the model and empirical Bayes disagree show where pitch-level information moves a projection away from a hitter's outcome history. Across 545 hitters in 2024, the two methods' projected splits have a Spearman correlation of 0.77, falling from 0.80 and 0.79 in the low and medium strata to 0.69 in the high stratum (Figure~\ref{fig:disagreement}). Lower agreement in the high stratum is consistent with empirical Bayes shrinking least for hitters with long outcome histories, while the model's projection still rests on how the hitter handles each pitch. Ranking within batting side, the largest disagreement is Robbie Grossman, first among right-handed hitters by empirical Bayes and 324th by the model. Marcell Ozuna is the reverse, 6th by the model and 316th by empirical Bayes, and Aaron Judge and José Altuve also rank much higher by the model. Figure~\ref{fig:cases} grades Ozuna, Judge, and Altuve, and six medium-stratum hitters from the twenty largest disagreements, against 2025. Most observed splits there have a standard error between 0.04 and 0.06, and a few near 0.1, at least the size of the gaps being graded, so the cases show how the two methods differ rather than which is right. The evidence on accuracy is the 539-hitter comparison in Section~\ref{sec:test}.

% Source: paper/draft2/make_figs.py cases().
\begin{figure}[!ht]
\centering
\includegraphics[width=0.78\linewidth]{fig_cases.png}
\caption{\toreview 2025 platoon splits for the five case-study hitters (top), the three high-exposure disagreements in Table~\ref{tab:cases} (middle), and the six medium-stratum hitters among the twenty largest 2024 disagreements (bottom). Each observed split has a 1.96 standard error bar beside the model and empirical Bayes projections. Every bar holds both projections.}
\label{fig:cases}
\end{figure}

\begin{figure}[!ht]
\centering
\includegraphics[width=0.6\linewidth]{fig_disagreement.png}
\caption{\toreview Rank of each hitter's projected 2024 platoon split within batting side, model (vertical) against empirical Bayes (horizontal), for 545 hitters colored by prior-exposure stratum. Rank 1 favors left-handed pitchers most. The dashed diagonal is equal rank, and the labeled hitters are large disagreements, including those named in the text.}
\label{fig:disagreement}
\end{figure}

\subsection{What the model sees in a hitter}
\label{sec:sees}

% Source: results/v_chain_clean/surfaces_location.csv (quantity q, delta x1000; selection model
% clean_clean_dim64); make_figs3.py fig_surfaces.
Because the network scores single pitches, it can be read one pitch location at a time. Figure~\ref{fig:surfaces} shows, for each of the five case-study hitters, the expected wOBA on a ball in play by location, relative to the average of 100 hitters on the same pitches. The surfaces differ in shape. Schwarber and Soto do their damage on pitches in the upper middle of the zone, up to about 170 points above average, and Soto falls below average on low pitches at and below the bottom of it. Trout's damage sits low, about 80 points above average along the bottom row and close to average at the top, peaking at about 200 to 220 points on low pitches over the middle. Bohm's surface is nearly flat, with an average edge of about 30 to 40 points and no cell above about 70. Each hitter's surface also changes little between pitcher hands. Pederson's average edge on contact is 44 points against left-handed pitchers and 48 against right-handed pitchers, so the quality of contact differs little between hands (Section~\ref{sec:anchors}).

\begin{figure}[!ht]
\centering
\includegraphics[width=\linewidth]{fig_surfaces.png}
\caption{\toreview Expected wOBA on a ball in play by pitch location, hitter minus the average of 100 hitters on the same 2015 to 2023 pitches, in wOBA points, against left-handed (top) and right-handed (bottom) pitchers, for Pederson, Schwarber, Soto, Trout, and Bohm (columns). Red is above the average hitter and blue below. Locations are from the catcher's view, with ``in'' toward the batter. The dotted box is the strike zone.}
\label{fig:surfaces}
\end{figure}

% Source: make_figs3.py fig_miss_damage (asserts match to v_chain_clean/embedding_structure/
% embedding_alt_views.json); decision log (spin vs fastball pairings).
The embedding separates some skills by pitch type and not others. Figure~\ref{fig:missdamage} fits two directions in the embedding for each of two skills, one for breaking balls and one for fastballs. For whiff rate the two directions correlate at 0.64 across hitters, and many hitters sit well off the diagonal, standing much worse against one pitch family than the other relative to other hitters. Their observed whiff rates on the two families correlate at 0.53, so the embedding keeps a distinction that is present in the data. For power, measured by 90th-percentile exit velocity, the two directions correlate at 0.94 and the observed values at 0.83, close to a single direction. One possible reason is that exit velocity is observed only on balls in play, which leaves fewer labels per pitch family than swings and contact have, but this was not tested. Both correlations in the embedding run higher than the observed ones, partly because the two directions share the space's dominant variation.

\begin{figure}[!ht]
\centering
\includegraphics[width=0.8\linewidth]{fig_miss_damage.png}
\caption{\toreview Each hitter placed on embedding directions fit to breaking-ball and fastball whiff rate (left, 1,335 hitters) and to breaking-ball and fastball power, 90th-percentile exit velocity (right, 1,203 hitters), in standard deviations. The dashed line is equal standing on both pitch families. The correlations of the observed hitter statistics are given for comparison.}
\label{fig:missdamage}
\end{figure}

% Source: results/v_chain_clean/neighbours_high.csv; embedding_structure.json (silhouette max 0.054, k 2 to 10).
Nearest neighbors in the embedding show what the model treats as similar. Among high-exposure hitters, Pederson's eight nearest neighbors are Kyle Schwarber, Joey Gallo, Justin Smoak, Harrison Bader, Logan Morrison, Franmil Reyes, Ryan Zimmerman, and Matt Chapman, mostly left- and right-handed power bats. Trout's are Ronald Acuña, Pete Alonso, Brandon Belt, Eugenio Suárez, Austin Jackson, Anthony Rendon, Howie Kendrick, and Christian Walker. Ranked by ten observed statistics instead, covering plate discipline, contact quality, and wOBA, Trout's nearest hitters are Yordan Álvarez, Paul Goldschmidt, Juan Soto, Mookie Betts, Joey Votto, Rhys Hoskins, Ronald Acuña, and Will Smith, and only Acuña appears on both lists. Across the five case-study hitters the two kinds of list share 9 of 40 names. The embedding's notion of similarity differs from one built on season statistics. The neighbor search also uses direction only, ignoring the vector length, which tracks playing time (correlation 0.93 with log prior PA). These neighborhoods shade into one another. The silhouette score, near zero when groups overlap, is at most 0.054 for any number of clusters from two to ten, so the space is better read as neighborhoods of similar hitters than as a set of hitter types.

\subsection{Other queries}
\label{sec:otherqueries}

% Source: results/v_chain_clean/paper_figures/level_query_heldout.json (partial on log prior PA,
% held-out 2024, min 10 PA); results/pitcher_type_query_clean/ceiling.csv (split-half reliability).
Any set of pitchers defines a query, and only the handedness query was tested against baselines. Two others illustrate what the approach supports and what limits checking it. The level query averages a hitter's projections against left- and right-handed pitchers, each weighted over every pitcher of that hand (Figure~\ref{fig:levelquery}). Its correlation with held-out 2024 wOBA, after removing log prior PA, is 0.38 across 521 hitters (interval 0.30 to 0.45), and it rises with exposure, from 0.17 in the low stratum to 0.38 in the medium and 0.52 in the high. A query over pitcher traits, such as right-handed pitchers in the top third of fastball velocity, runs the same way. Checking it is harder than checking the projection. A hitter's observed wOBA against the top third of right-handed pitchers by fastball velocity, minus that against the bottom third, has a split-half reliability of 0.08 across 301 hitters in 2024 (interval $-0.15$ to 0.26), and the same right-handed contrasts for fastball height and breaking-ball share are no higher. One season of outcomes cannot distinguish these matchup effects from noise, so the model's projections for them remain untested.

\begin{figure}[!ht]
\centering
\includegraphics[width=0.9\linewidth]{fig_level_query.png}
\caption{\toreview The level query against held-out 2024 wOBA for 521 hitters with at least 10 PA. Left, each hitter's projection against all pitchers and observed wOBA, colored by prior-exposure stratum, with point size proportional to 2024 PA. Right, the partial correlation after removing log prior PA, with 95\% intervals.}
\label{fig:levelquery}
\end{figure}


